% Canonical authority: composition, responses, covariance and frozen quotation.
\newcommand{\CompositionPage}{%
\pageunit{Removal and rejoin determine the conditional result}{sec:composition}{%
Bind $\Pi_k:M_k\to Y_k$, $F_k:Y_k\to Z_k$ and $B_k:M_k\to Z_k$.
$\Pi_k(\ma)$ is the declared removal contribution; $B_k(\ma)$ is the
declared contribution at rejoin. Only when the two operand pairs have the
compatibility in Section~\ref{sec:types} does the ordinary relation apply:
\begin{equation}\begin{aligned}
\rho_k&=\yin-\Pi_k(\ma),\\
z_k&=F_k(\rho_k)+B_k(\ma).
\end{aligned}\label{eq:composition}\end{equation}
The model path bypasses the exactly named residual operation. The operation
graph binds order, removal/rejoin locations, all path application counts and
consequential state. It does not imply bypass of every upstream operation.
There is no assumed $\Pi_k=B_k$, inverse relation, equal domain/support,
linear $F_k$, true model or unit response. Missing path contributions are not
zero and undefined partial-domain maps are not silently extended.

\subsection*{What follows when only the residual operation is linear}
If the complete $F_k=A_k$ is exact fixed linear on the declared spaces,
support, state and conditioning, distributivity gives
\begin{equation}
z_k=A_k\yin-A_k[\Pi_k(\ma)]+B_k(\ma).
\label{eq:mapvalued}\end{equation}
This remains a map-valued expression. Nonlinear $\Pi_k$ or $B_k$ prevents a
linear applied-model operator from being inferred. A fixed loading basis
alone is insufficient; coefficient recomputation alone neither proves nor
disproves complete linearity. Centering, support and every material state
belong to the claim. An affine offset cannot be discarded to obtain this form.
Other procedures retain Equation~\ref{eq:composition} and their full state graph.

\subsection*{The stronger, completely fixed linear specialization}
Only if $A_k,\Pi_k,B_k$ are \emph{all} exact fixed linear maps on the stated
spaces, units, support, state and conditioning may one define
\begin{equation}
D_k:=B_k-A_k\circ\Pi_k:M_k\to Z_k,\qquad
z_k=A_k\yin+D_k\ma.
\label{eq:linear}\end{equation}
The difference operator is well typed because both terms map the applied-model
space into one compatible result space. It describes the conditional net
influence of applied feedback, not all sky bias. Holding the same input and
operators and changing $\ma=m_{\rm ref}^{\ap}+\delta_m^{\ap}$ gives
\begin{equation}
z_k(\ma)-z_k(m_{\rm ref}^{\ap})=D_k\delta_m^{\ap}.
\label{eq:modelerror}\end{equation}
Its conditional mean is the corresponding model-induced bias component.
Accepted-model error requires its declared transformation first. Zero/identity
limiting cases in the predictions are deductions under premises, not defaults
or admitted run plans. None restores an upstream unavailable measured mode.
\src{RESPONSE\_AND\_UNCERTAINTY\_FAMILIES; SCOPE\_BRIEF \S2; REQUIRED\_CONDITIONAL\_PREDICTIONS.}
}}
\newcommand{\ResponseLocalPage}{%
\pageunit{Response depends on what is perturbed and what is held fixed}{sec:response-local}{%
A response is an exact query about a procedure. Bind perturbation source/domain,
direction/basis, units, limit or finite step and normalization when relevant;
fixed versus rerun state; recurrence and terminal scope; baseline/perturbed
comparability; branch/support/state changes; codomain; derivative or
finite-difference convention; application counts, generations, conditioning;
and unavailable/discontinuity behavior. Reference/gauge/null-space comparisons
must be compatible. A changed reference is not silently a signal perturbation.
No numerical query or step is selected here.

\begin{description}[style=nextline,leftmargin=0pt,labelwidth=0pt,labelsep=0pt]
\item[RF-01: fixed one-step measured-input response]
Perturb $\yin$ in $Y_k$, holding the applied model and declared maps, state
and support fixed. Evaluate the unchanged one-step composition; no learning,
recurrence, stopping or terminal selection. Codomain: $Z_k$ differences.
Measured-input names this path, not necessarily the original parent or an
independent observation.
\item[RF-02: fixed one-step applied-model response]
Perturb $\ma$ in $M_k^{\ap}$, holding the measured iteration input and
declared maps, state and support fixed. Evaluate the unchanged one-step
composition; no acceptance/model learning, recurrence or selection.
Codomain: $Z_k$ differences. Candidate or accepted perturbations require their
transformations in the corresponding full procedure.
\item[RF-03: one-iteration FRUIT full-procedure response]
Perturb the declared one-iteration input on its exact domain. Hold prior $S_k$
and upstream producer procedure/state fixed. Rerun construction, selection,
state resolution, support, coefficient realization and application for
iteration $k$ only. Codomain: its $Z_k$ response. No earlier recurrence or
terminal selector belongs to this local role.
\item[RF-04: fixed-sequence recursive response]
Perturb the declared recursion input. Hold method rules, initial state,
upstream producer procedure/state and the predeclared fixed sequence or
absolute endpoint index fixed. Propagate input/model/state dependence through
that extent, binding exact input and result/sequence codomains. Data-dependent
sequence length and terminal selection are outside this role.
\end{description}
All roles retain branch/support/state changes and typed unavailability or
transitions when comparisons/results/derivatives do not exist. Under complete
fixed linearity in Equation~\ref{eq:linear}, the separate one-step responses are
\begin{equation}
R_{z\leftarrow y,k}^{\mathrm{FRUIT,fixed}}=A_k,\qquad
R_{z\leftarrow m^{\ap},k}^{\mathrm{FRUIT,fixed}}=D_k.
\label{eq:responses}\end{equation}
If only $F_k=A_k$ is linear, the fixed input response may still be $A_k$;
the model response needs the derivative of the complete map-valued expression
where it exists, not an inferred $D_k$.
\src{RESPONSE\_AND\_UNCERTAINTY\_FAMILIES; NOTATION\_AND\_ROLE\_TAXONOMY.}
}}
\newcommand{\ResponseOuterPage}{%
\pageunit{Terminal, FRUIT-only and whole-chain queries have different reach}{sec:response-outer}{%
\begin{description}[style=nextline,leftmargin=0pt,labelwidth=0pt,labelsep=0pt]
\item[RF-05: selected-terminal response]
Perturb the declared selection/procedure input and name all fixed initial and
upstream facts. Include consequential iterations, convergence, stopping,
resource limits and terminal selection, with complete candidate population and
causes. Bind input and selected-result codomain, query convention and terminal
record. Include branch/support/state/selection transitions and reference/gauge/
null-space comparison. Unavailable dependencies or nonexistent derivatives
remain unavailable or typed transitions.
\item[RF-06: whole-chain response]
Perturb the exact upstream/source query on its domain, naming fixed external
facts. Rerun every included data-dependent upstream producer and the declared
FRUIT procedure, including its recurrence/terminal scope. Bind source and
final domains and convention; each included producer and comparison must be
available. Otherwise retain exact unavailability/discontinuity.
\item[RF-07: complete FRUIT-only parent-domain response]
Perturb the realized parent-domain payload while holding the upstream producer
operation, learned state and generation fixed. Compose the parent-domain
entry with the complete FRUIT procedure, including recurrence and terminal
selection as applicable. Bind parent-query/final-response domains, initial
FRUIT conditions, query convention and RF-03/RF-04/RF-05 component references.
Include branch/support/state/selection and reference/gauge/null-space
consequences. No upstream producer is rerun.
\end{description}
RF-03 supplies local iteration relations. RF-04 propagates fixed extent;
RF-05 includes the selector and data-dependent stopping extent. RF-07 binds
the outer realized-parent query to those components: it is not an independent
alias of RF-03, RF-04 or RF-05. When a one-iteration special case reduces to a
component, record that exact specialization and reuse its bound response/status.
Equal values do not erase query identity. RF-06 additionally reruns included
upstream producers. The previous RF-07 held-numerical-parent meaning is retired.

Perturbed parent payloads are separately identified response-query objects
with immutable baseline-parent ancestry. A fixed producer generation does
not mutate the measured parent or claim a new upstream realization.
A rule-changing control perturbation requires a new method identity and
explicit comparison of methods, not a state perturbation disguised as one rule.

For RF-01/RF-02, ``fixed'' means the conditioned map and unchanged application
rule. Coefficients recomputed by that rule from a perturbed operand must still
be recomputed, with new coefficient/application states. Freezing their numerical
values instead changes the procedure. Missing baseline/perturbed codomain
compatibility makes the response unavailable. Threshold, branch, support,
state or selector changes may defeat derivative existence; a finite difference
needs its own convention, never a substituted fixed-branch derivative.

Use all response identities/statuses required by the applicable method and
claim. The seven-role catalog does not make any numerical response universally
available. No online adaptive estimator is introduced by these names.
\src{RESPONSE\_AND\_UNCERTAINTY\_FAMILIES; SCIENTIFIC\_OWNER\_DECISION\_LEDGER (response crosswalk).}
}}
\newcommand{\UncertaintyPage}{%
\pageunit{Shared input and model information must remain in uncertainty}{sec:uncertainty}{%
Assume complete fixed linearity, compatible real vector representations,
finite joint second moments and exact conditioning $\Om$ fixing relevant
operators/state, reference/gauge/null-space coordinates and support. Define
$C_{yy}^{\mathrm{FRUIT}},C_{mm}^{\mathrm{FRUIT}},C_{ym}^{\mathrm{FRUIT}},
C_{my}^{\mathrm{FRUIT}}$ as the joint conditional covariance blocks of
$(\yin,\ma)$, with exact quantity axes, units, domains and conditioning.
The $m$ axis denotes only the applied model. Equation~\ref{eq:linear} then gives
\begin{equation}\begin{aligned}
C_z^{\mathrm{FRUIT}}={}&A_k C_{yy}^{\mathrm{FRUIT}} A_k^{\mathsf T}
 +D_k C_{mm}^{\mathrm{FRUIT}}D_k^{\mathsf T}\\
 &+A_k C_{ym}^{\mathrm{FRUIT}}D_k^{\mathsf T}
 +D_k C_{my}^{\mathrm{FRUIT}}A_k^{\mathsf T}.
\end{aligned}\label{eq:covariance}\end{equation}
This follows by expanding the conditional covariance of the sum; the two
mixed products are the two cross terms. In a valid real joint covariance,
$C_{my}^{\mathrm{FRUIT}}=(C_{ym}^{\mathrm{FRUIT}})^{\mathsf T}$.
Do not drop either term without exact independence/conditioning authority
establishing its absence. A model learned from the parent is not independent
by implication. Candidate/accepted covariance needs the exact transformation
and its dependence before it can supply applied covariance.

If the applied model is deterministic under the conditioning, its variance
and cross blocks vanish there, leaving $A_k C_{yy}^{\mathrm{FRUIT}}A_k^{\mathsf T}$.
This conditional target omits model-learning variation; it is not unconditional
independence or covariance completeness and cannot erase unknown sources from
another target. Random-operator/nonlinear procedures need their own uncertainty
method. A Jacobian approximation is separately typed, not this exact identity.

Keep nine uncertainty roles distinct: measured-parent covariance;
candidate/accepted/applied model/prior covariance; parent-model cross
covariance; operator/state uncertainty; support/threshold selection;
stopping/terminal selection; model mismatch/bias; empirical repeatability;
and NOI uncertainty. Bias is not covariance by identity. Each binds reference/
gauge, unavailable modes, model/prior support, omitted mode/reference variation
and exact conditioning. A fixed gauge or singular covariance does not mean an
unavailable absolute mode was measured with zero uncertainty. Selected-terminal
uncertainty additionally binds the full candidate population, stop/resource
facts, selection record and causes. Reused parents imply no independent
observations or $1/\sqrt{\text{iteration count}}$ uncertainty law.

Fixed-state NOI applies identical fixed FRUIT state/operator graph with exact
application parity to every compatible admitted realization; absent parity
makes the numerical route unavailable. It omits unmodeled learning, support,
selection, convergence, restart and leakage variation. NOI-informed later
learning is a dependent immutable successor science/ensemble/uncertainty
generation, not independent validation or a retroactive edit. Changing to
per-realization relearning of model, support, cleaner, coefficients, stopping
or selector requires a separate method/ensemble and complete member-state graph.
Members under unchanged rules get new state/application generations, not a
method each. Fixed/relearned members cannot share one ensemble or be pooled here.
\src{RESPONSE\_AND\_UNCERTAINTY\_FAMILIES; BOUNDARIES\_AND\_CONVENTIONS B4.}
}}
\newcommand{\PTCAppendix}{%
\pageunit{Frozen PTC application quotation}{app:ptc}{%
This is the byte-exact permitted SCI-PTC v0.1/r0.5 application fragment,
under the source-local notation cover in Appendix~\ref{app:notation}.
It supplies no numerical FRUIT cleaning, rank, metric, model or parent route.
FRUIT imports the distinction and restrictions, not a fresh PTC derivation.
\begin{quote}\smallFor the ordinary configured-rank PCA/SVD route, application uses the frozen
group-local detector-loading subspace and mask-aware linear coefficient
recomputation. For centered row
$x_{g,t}=Y^{\rm CAL}_{g,t}-\lambda_g$,
\begin{equation}
\begin{aligned}
G^{\rm app}_{g,t}
&:=\widehat A_g^{\mathsf T}W^{\rm app}_{g,t}\widehat A_g,
&\operatorname{rank}_{\tau_g}(G^{\rm app}_{g,t})
&=k_{{\rm req},g},\\
\widehat m^{\rm app}_{g,t}
&=x_{g,t}W^{\rm app}_{g,t}\widehat A_g
  \left(G^{\rm app}_{g,t}\right)^+,\\
z_{g,t}
&=x_{g,t}-\widehat m^{\rm app}_{g,t}\widehat A_g^{\mathsf T}.
\end{aligned}
\label{eq:conditional-cleaner}
\end{equation}
The exact output support determines which entries of $z_{g,t}$ are published.
The solve is authorized only when every required input and
$G^{\rm app}_{g,t}$ is finite and its numerical rank under frozen tolerance
$\tau_g$ equals $k_{{\rm req},g}$. Otherwise application and output for that
group-time are unavailable with cause. The generalized inverse does not
authorize partial-rank subtraction, rank reduction, interpolation, or
borrowing detectors from another group. Detector binding, metric, support,
and insufficient-support state are frozen in $\Theta_g$. On complete uniform
support this reduces to detector-right projection through
$I-\widehat A_g(\widehat A_g^{\mathsf T}\widehat A_g)^+
\widehat A_g^{\mathsf T}$. Any temporal-left, two-sided,
detector/time-specific, or general vectorized family remains separately named
and is not inferred from matrix orientation or rank.

A separately named frozen-component subtraction family is
\begin{equation}
\widehat U_{\rm total}:=U_{{\rm total},\Theta}(Y^{\rm CAL}),
\qquad
\mathcal A^{\rm sub}_{\widehat U_{\rm total}}(Y')=Y'-\widehat U_{\rm total},
\qquad
D\mathcal A^{\rm sub}_{\widehat U_{\rm total}}[Y](H)=H.
\label{eq:frozen-component}
\end{equation}
It is not interchangeable with Equation~\ref{eq:conditional-cleaner} and is
not the base-v0.1 PCA/SVD application rule.
\end{quote}
\src{PTC\_APPLICATION\_REFERENCE.tex, exactly preserved; BOUNDARIES\_AND\_CONVENTIONS B2.}
}}
